\chapter{Spectral Theory}

Throughout this chapter, scalar fields are complex whenever the spectrum of an
operator is discussed.

\section{Closed operators, resolvents, and spectra}

\begin{notedefinition}
Let $X$ be a complex Banach space and let
$T\colon\Dom(T)\subseteq X\to X$ be linear.  A number $\lambda\in\mathbb C$
is a \emph{regular value} if
\[
  T-\lambda I\colon\Dom(T)\to X
\]
is bijective and its inverse
\[
  R_\lambda(T)=(T-\lambda I)^{-1}\colon X\to X
\]
is bounded.  The \emph{resolvent set} $\rho(T)$ is the set of regular values,
and the \emph{spectrum} is
\[
  \sigma(T)=\mathbb C\setminus\rho(T).
\]
The spectrum is partitioned into:
\begin{itemize}
  \item the point spectrum, where $T-\lambda I$ is not injective;
  \item the continuous spectrum, where it is injective and has dense but
        non-surjective range;
  \item the residual spectrum, where it is injective but its range is not
        dense.
\end{itemize}
\end{notedefinition}

\begin{noteproposition}
If $T$ is closed and $T-\lambda I$ is bijective, then
$R_\lambda(T)$ is automatically bounded.  In particular, for a closed
operator the definition of $\rho(T)$ may omit boundedness of the inverse.
\end{noteproposition}

\begin{proof}
The inverse has graph obtained from $\Graph(T-\lambda I)$ by interchanging the
two coordinates, so it is closed.  Its domain is all of the Banach space $X$;
the closed graph theorem makes it bounded.
\end{proof}

\begin{notetheorem}[Neumann series]
Let $A\in\Bop{X}{X}$ with $\|A\|<1$.  Then $I-A$ is invertible and
\[
  (I-A)^{-1}=\sum_{n=0}^{\infty}A^n
\]
in operator norm.  More generally, if $\lambda_0\in\rho(T)$ and
$|\lambda-\lambda_0|\,\|R_{\lambda_0}(T)\|<1$, then $\lambda\in\rho(T)$ and
\[
  R_\lambda(T)
  =R_{\lambda_0}(T)
   \sum_{n=0}^{\infty}
     \bigl((\lambda-\lambda_0)R_{\lambda_0}(T)\bigr)^n.
\]
\end{notetheorem}

\begin{proof}
The partial sums telescope:
$(I-A)\sum_{n=0}^N A^n=I-A^{N+1}$, and $A^{N+1}\to0$.  For the resolvent,
factor
\[
  T-\lambda I
  =\bigl(I-(\lambda-\lambda_0)R_{\lambda_0}(T)\bigr)
    (T-\lambda_0I).
\]
\end{proof}

\begin{notecorollary}
The resolvent set is open and the spectrum is closed.  If $T\in\Bop{X}{X}$,
then
\[
  \sigma(T)\subseteq\{\lambda:|\lambda|\le\|T\|\}.
\]
\end{notecorollary}

\begin{proof}
Openness follows from the local Neumann-series formula.  If
$|\lambda|>\|T\|$, then
\[
  T-\lambda I=-\lambda\left(I-\frac{T}{\lambda}\right)
\]
is invertible by the Neumann series.
\end{proof}

\begin{notetheorem}[Resolvent identity]
For $\lambda,\mu\in\rho(T)$,
\[
  R_\lambda(T)-R_\mu(T)
  =(\lambda-\mu)R_\lambda(T)R_\mu(T).
\]
In particular, the resolvents commute.
\end{notetheorem}

\begin{proof}
Use $A^{-1}-B^{-1}=A^{-1}(B-A)B^{-1}$ with
$A=T-\lambda I$ and $B=T-\mu I$.
\end{proof}

\begin{notetheorem}
For every bounded operator on a nonzero complex Banach space, the spectrum is
nonempty and compact.
\end{notetheorem}

\begin{proof}
Closedness and boundedness were proved above.  If $\sigma(T)=\varnothing$,
then for each $x\in X$ and $f\in X^*$ the scalar function
\[
  \lambda\longmapsto f(R_\lambda(T)x)
\]
is entire.  For $|\lambda|>\|T\|$, the Neumann expansion gives
$\|R_\lambda(T)\|\le(|\lambda|-\|T\|)^{-1}$, so this entire function tends to
zero at infinity.  Liouville's theorem makes it identically zero.  Since
functionals separate points, $R_\lambda(T)=0$, impossible for an inverse.
\end{proof}

\begin{notedefinition}
For $T\in\Bop{X}{X}$, the \emph{spectral radius} is
\[
  r(T)=\sup\{|\lambda|:\lambda\in\sigma(T)\}.
\]
\end{notedefinition}

\begin{noteproposition}
Let $T\in\Bop{X}{X}$.  If $\lambda\in\rho(T)$, then $R_\lambda(T)$ commutes
with $T$.  More generally,
$R_\lambda(T)$ commutes with every bounded operator that commutes with $T$.
Consequently, all resolvents commute with one another.
\end{noteproposition}

\begin{proof}
From
\[
  (T-\lambda I)R_\lambda(T)
  =R_\lambda(T)(T-\lambda I)=I
\]
we obtain
\[
  TR_\lambda(T)=I+\lambda R_\lambda(T)
  =R_\lambda(T)T.
\]
Now let $S\in\Bop{X}{X}$ commute with $T$.  Then $S$ commutes with
$T-\lambda I$, and hence
\[
  (T-\lambda I)SR_\lambda(T)
  =S(T-\lambda I)R_\lambda(T)=S.
\]
Applying $R_\lambda(T)$ on the left gives
$SR_\lambda(T)=R_\lambda(T)S$.  Taking
$S=R_\mu(T)$ proves that any two resolvents commute.
\end{proof}

\section{Spectral mapping and the spectral radius formula}

\begin{notetheorem}[Polynomial spectral mapping theorem]
Let $T\in\Bop{X}{X}$ and let $p$ be a complex polynomial.  Then
\[
  \sigma(p(T))=p(\sigma(T)).
\]
\end{notetheorem}

\begin{proof}
First let $\mu\notin p(\sigma(T))$.  Factor
\[
  p(z)-\mu=c\prod_{j=1}^m(z-\lambda_j).
\]
Every $\lambda_j$ lies in $\rho(T)$, so
$p(T)-\mu I=c\prod_j(T-\lambda_jI)$ is invertible.

Conversely, let $\mu=p(\lambda)$ with $\lambda\in\sigma(T)$.  Polynomial
division gives
\[
  p(T)-p(\lambda)I=(T-\lambda I)q(T),
\]
and the two factors commute.  If the left side were invertible, then
$(T-\lambda I)$ would have both a bounded left inverse and a bounded right
inverse obtained by combining $q(T)$ with that inverse, hence would itself be
invertible.  This contradicts $\lambda\in\sigma(T)$.
\end{proof}

\begin{noteproposition}[Linear independence of eigenvectors]
Eigenvectors corresponding to distinct eigenvalues are linearly independent.
\end{noteproposition}

\begin{proof}
If $\sum_{j=1}^n a_jx_j=0$ is a shortest nontrivial relation and
$Tx_j=\lambda_jx_j$, apply $T-\lambda_nI$.  The resulting relation among
$x_1,\ldots,x_{n-1}$ has coefficients
$a_j(\lambda_j-\lambda_n)$ and is nontrivial, a contradiction.
\end{proof}

\begin{notetheorem}[Spectral radius formula]
For $T\in\Bop{X}{X}$,
\[
  r(T)=\lim_{n\to\infty}\|T^n\|^{1/n}
      =\inf_{n\ge1}\|T^n\|^{1/n}.
\]
\end{notetheorem}

\begin{proof}
By spectral mapping,
$r(T)^n=r(T^n)\le\|T^n\|$, so
$r(T)\le\inf_n\|T^n\|^{1/n}$.

Let $R>r(T)$.  The resolvent is analytic on $|\lambda|>r(T)$ and has the
Laurent expansion at infinity
\[
  R_\lambda(T)
  =-\sum_{n=0}^{\infty}\frac{T^n}{\lambda^{n+1}}
\]
for $|\lambda|>\|T\|$.  Cauchy's formula, applied on the circle
$|\lambda|=R$, gives
\[
  T^n=-\frac{1}{2\pi i}\int_{|\lambda|=R}
          \lambda^nR_\lambda(T)\,\dd\lambda.
\]
Hence $\|T^n\|\le C_RR^{n+1}$, and
$\limsup_n\|T^n\|^{1/n}\le R$.  Letting $R\downarrow r(T)$ gives the reverse
inequality.  The equality with the infimum follows from submultiplicativity of
$\|T^n\|$ and Fekete's lemma applied to $\log\|T^n\|$.
\end{proof}

\section{Compact operators}

\begin{notedefinition}
A linear map $T\colon X\to Y$ is compact if it maps bounded sets to relatively
compact sets.  Equivalently, for every bounded sequence
$(x_n)$ in $X$, the sequence $(Tx_n)$ has a convergent subsequence in $Y$.
\end{notedefinition}

\begin{notelemma}
Every compact linear operator is bounded.  The identity operator on a normed
space is compact if and only if the space is finite dimensional.
\end{notelemma}

\begin{proof}
If $T$ were unbounded, choose $x_n$ with $\|x_n\|\le1$ and $\|Tx_n\|\ge n$;
then $(Tx_n)$ has no convergent subsequence.  The statement about the identity
is the Riesz lemma characterization of finite-dimensional normed spaces.
\end{proof}

\begin{notetheorem}[Sequential criterion]
A linear map $T\colon X\to Y$ is compact if and only if every bounded sequence
$(x_n)$ has a subsequence $(x_{n_k})$ for which $(Tx_{n_k})$ converges.
\end{notetheorem}

\begin{proof}
If $T$ maps bounded sets to relatively compact sets and $(x_n)$ is bounded,
then all $Tx_n$ lie in a set with compact closure, so a subsequence converges.

Conversely, let $B\subseteq X$ be bounded.  Every sequence in $T(B)$ has the
form $(Tx_n)$ with $x_n\in B$, and therefore has a convergent subsequence.
Thus $T(B)$ is relatively sequentially compact.  To pass to the closure, let
$(y_n)\subseteq\overline{T(B)}$ and choose $z_n\in T(B)$ with
$\|y_n-z_n\|<1/n$.  A subsequence of $(z_n)$ converges, and the corresponding
subsequence of $(y_n)$ has the same limit.  Hence $\overline{T(B)}$ is
sequentially compact.  In a metric space, sequential compactness and
compactness are equivalent, so $\overline{T(B)}$ is compact and $T$ is
compact.
\end{proof}

\begin{notetheorem}
Every bounded finite-rank operator is compact.  If $Y$ is Banach, then
$\Kop{X}{Y}$ is closed in $\Bop{X}{Y}$ for the operator norm.
\end{notetheorem}

\begin{proof}
A bounded set in a finite-dimensional range is relatively compact.  For the
closure statement, let $T_n\to T$ in operator norm with each $T_n$ compact.
Let $(x_j)$ be bounded, say $\|x_j\|\le M$.  Compactness of $T_1$ gives a
subsequence on which $T_1x_j$ converges.  From that subsequence choose one on
which $T_2x_j$ converges, and continue.  If $(y_j)$ is the diagonal
subsequence, then for every fixed $n$ the sequence $(T_ny_j)$ converges.

Fix $\varepsilon>0$ and choose $n$ so large that
\[
  2M\|T-T_n\|<\frac{2\varepsilon}{3}.
\]
For sufficiently large $j,k$,
\(
  \|T_ny_j-T_ny_k\|<\varepsilon/3
\), and hence
\[
\begin{aligned}
  \|Ty_j-Ty_k\|
   &\le \|(T-T_n)y_j\|
       +\|T_ny_j-T_ny_k\|
       +\|(T_n-T)y_k\|\\
   &<\varepsilon.
\end{aligned}
\]
Thus $(Ty_j)$ is Cauchy and converges because $Y$ is Banach.  The sequential
criterion proves that $T$ is compact.
\end{proof}

\begin{notedefinition}
A metric-space subset $B$ is \emph{totally bounded} if, for every
$\varepsilon>0$, it is covered by finitely many balls of radius
$\varepsilon$.
\end{notedefinition}

\begin{notelemma}
A metric space is compact if and only if it is complete and totally bounded.
Every totally bounded set is separable, and every subset of a totally bounded
set is totally bounded.
\end{notelemma}

\begin{proof}
Suppose first that $M$ is compact.  Compact metric spaces are complete.  For
each $\varepsilon>0$, the open cover
\(
  \{B(x,\varepsilon):x\in M\}
\)
has a finite subcover, so $M$ is totally bounded.

Conversely, assume that $M$ is complete and totally bounded, and let
$(x_n)$ be a sequence in $M$.  Cover $M$ by finitely many balls of radius
$2^{-1}$.  One of them contains infinitely many terms; retain that
subsequence.  Cover $M$ by finitely many balls of radius $2^{-2}$ and retain
an infinite subsequence lying in one of them.  Continuing in this way and
taking the diagonal subsequence $(x_{n_k})$, we have
\[
  d(x_{n_j},x_{n_k})\le 2^{1-m}
  \qquad(j,k\ge m),
\]
because both points lie in the same ball of radius $2^{-m}$.  Thus the
diagonal subsequence is Cauchy and converges in $M$ by completeness.  Hence
$M$ is sequentially compact, and therefore compact.

If $B$ is totally bounded, choose for each $n$ a finite $1/n$-net $F_n$.
The countable set $\bigcup_nF_n$ is dense in $B$, so $B$ is separable.  If
$A\subseteq B$, a finite $\varepsilon/2$-net for $B$ gives a finite
$\varepsilon$-net for $A$ after discarding balls that miss $A$ and choosing
one point of $A$ from each remaining ball.  Thus every subset of a totally
bounded set is totally bounded.
\end{proof}

\begin{notetheorem}
If $T\colon X\to Y$ is compact, then $\Ran(T)$ is separable.
\end{notetheorem}

\begin{proof}
Write
\[
  \Ran(T)=\bigcup_{n=1}^{\infty}T(nB_X).
\]
Each $T(nB_X)$ is totally bounded and therefore separable.  A countable union
of countable dense sets is dense in the range.
\end{proof}

\begin{notetheorem}[Compact extension from a dense subspace]
Let $X$ and $Y$ be Banach spaces, let $D\subseteq X$ be dense, and let
$T\colon D\to Y$ be bounded.  Its unique bounded extension $\widetilde T\colon X\to Y$ is compact
if $T$ maps the unit ball of $D$ to a totally bounded subset of $Y$.
\end{notetheorem}

\begin{proof}
First construct the extension.  If $x\in X$, choose $d_k\in D$ with
$d_k\to x$.  Since $T$ is bounded,
\[
  \|Td_k-Td_\ell\|\le\|T\|\,\|d_k-d_\ell\|,
\]
so $(Td_k)$ is Cauchy and converges in the Banach space $Y$.  Define
$\widetilde Tx=\lim_kTd_k$.  The estimate above shows that this value is
independent of the approximating sequence, and it also gives linearity,
$\|\widetilde T\|=\|T\|$, and uniqueness.

Now let $(x_n)$ be bounded, say $\|x_n\|\le M$.  Choose $d_n\in D$ with
\[
  \|x_n-d_n\|<\frac1n.
\]
Then $\|d_n\|\le M+1$ for all large $n$.  The set
$T((M+1)B_D)$ is totally bounded because it is a scalar multiple of
$T(B_D)$.  Hence $(Td_n)$ has a Cauchy subsequence $(Td_{n_k})$.  Moreover,
\[
  \|\widetilde Tx_{n_k}-Td_{n_k}\|
  \le\|T\|\,\|x_{n_k}-d_{n_k}\|\longrightarrow0.
\]
Thus $(\widetilde Tx_{n_k})$ is Cauchy and converges in $Y$.  By the
sequential criterion, $\widetilde T$ is compact.
\end{proof}

\begin{notetheorem}[Schauder theorem]
A bounded linear operator $T\colon X\to Y$ is compact if and only if
$T^*\colon Y^*\to X^*$ is compact.
\end{notetheorem}

\begin{proof}
Assume first that $T$ is compact.  Fix $\varepsilon>0$.  Since
$T(B_X)$ is totally bounded, choose $x_1,\ldots,x_m\in B_X$ such that for
every $x\in B_X$ there is $j$ with
\[
  \|Tx-Tx_j\|<\varepsilon.
\]
Consider the finite-dimensional coordinate map
\[
  A\colon Y^*\to\mathbb K^m,
  \qquad A(y^*)=(y^*(Tx_1),\ldots,y^*(Tx_m)).
\]
The set $A(B_{Y^*})$ is bounded in $\mathbb K^m$, hence totally bounded.
Choose $y_1^*,\ldots,y_N^*\in B_{Y^*}$ so that their coordinate vectors form
an $\varepsilon$-net for $A(B_{Y^*})$ in the supremum norm.

Given $y^*\in B_{Y^*}$, choose $k$ with
\[
  |(y^*-y_k^*)(Tx_j)|<\varepsilon
  \qquad(1\le j\le m).
\]
For $x\in B_X$, choose $j$ as above.  Since
$\|y^*-y_k^*\|\le2$,
\[
\begin{aligned}
 |(T^*y^*-T^*y_k^*)(x)|
  &=|(y^*-y_k^*)(Tx)|\\
  &\le |(y^*-y_k^*)(Tx_j)|
       +2\|Tx-Tx_j\|\\
  &<3\varepsilon.
\end{aligned}
\]
Taking the supremum over $x\in B_X$ gives
$\|T^*y^*-T^*y_k^*\|<3\varepsilon$.  Thus $T^*(B_{Y^*})$ is totally bounded;
its closure is complete in the Banach space $X^*$, hence compact.  Therefore
$T^*$ is compact.

Conversely, if $T^*$ is compact, the implication just proved shows that
$T^{**}$ is compact.  The canonical embeddings satisfy
\[
  T^{**}J_X=J_YT.
\]
If $(x_n)$ is bounded in $X$, compactness of $T^{**}$ gives a subsequence for
which $J_YTx_{n_k}$ converges in $Y^{**}$.  Because $J_Y(Y)$ is closed and
$J_Y$ is an isometry, $(Tx_{n_k})$ converges in $Y$.  Hence $T$ is compact.
\end{proof}

\begin{notelemma}[Ideal property]
For normed spaces $X,Y$, compact operators form a two-sided operator ideal:
if $T\in\Kop{X}{Y}$ and $A,B$ are bounded with compatible domains, then
$ATB$ is compact.
\end{notelemma}

\begin{proof}
Let $(x_n)$ be bounded.  Since $B$ is bounded, $(Bx_n)$ is bounded; compactness
of $T$ gives a subsequence for which $(TBx_{n_k})$ converges.  Continuity of
$A$ then makes $(ATBx_{n_k})$ converge.  The sequential criterion proves that
$ATB$ is compact.
\end{proof}

\section{The spectrum of a compact operator}

\begin{notetheorem}[Riesz spectral theorem for compact operators]
Let $T$ be a compact operator on an infinite-dimensional complex Banach space.
Then:
\begin{enumerate}
  \item every nonzero point of $\sigma(T)$ is an eigenvalue;
  \item for each $\lambda\ne0$, the eigenspace
        $\Ker(T-\lambda I)$ is finite dimensional;
  \item $\sigma(T)$ is at most countable, and $0$ is its only possible
        accumulation point.
\end{enumerate}
\end{notetheorem}

\begin{proof}
For $\lambda\ne0$, put $N_\lambda=\Ker(T-\lambda I)$.  On this subspace,
$T=\lambda I$, so
\[
  I_{N_\lambda}=\lambda^{-1}T|_{N_\lambda}
\]
is compact.  The identity on a normed space is compact only in finite
dimension; hence $N_\lambda$ is finite dimensional.

Fix $\varepsilon>0$ and suppose that there are infinitely many distinct
eigenvalues $\lambda_n$ with $|\lambda_n|\ge\varepsilon$.  Let
\[
  M_n=\Span\{v_1,\ldots,v_n\},
  \qquad Tv_j=\lambda_jv_j,
\]
where $v_j\ne0$.  Eigenvectors for distinct eigenvalues are linearly
independent, so $M_{n-1}\subsetneq M_n$.  Riesz's lemma gives
$x_n\in M_n$ with
\[
  \|x_n\|=1,
  \qquad \dist(x_n,M_{n-1})>\frac12.
\]
The subspaces $M_n$ are $T$-invariant, and
\[
  (T-\lambda_nI)x_n\in M_{n-1}.
\]
For $m<n$, also $Tx_m\in M_m\subseteq M_{n-1}$.  Therefore
\[
  Tx_n-Tx_m=\lambda_nx_n-z_{m,n}
  \qquad\text{for some }z_{m,n}\in M_{n-1},
\]
and hence
\[
  \|Tx_n-Tx_m\|
  =|\lambda_n|\left\|x_n-\frac{z_{m,n}}{\lambda_n}\right\|
  >\frac{\varepsilon}{2}.
\]
Thus $(Tx_n)$ has no Cauchy subsequence, contradicting compactness of $T$.
Only finitely many eigenvalues can therefore satisfy $|\lambda|\ge\varepsilon$.
Taking $\varepsilon=1/k$ proves countability of the nonzero eigenvalues and
shows that only $0$ can be an accumulation point.

The remaining assertion---that every nonzero spectral value is an
eigenvalue---follows from the Fredholm alternative proved below.
\end{proof}

\begin{noteproposition}
Let $T$ be compact and $\lambda\ne0$.  Then $\Ran(T-\lambda I)$ is closed.
\end{noteproposition}

\begin{proof}
Let $N=\Ker(T-\lambda I)$, which is finite dimensional, and choose a closed
complement $X=N\oplus M$.  The restriction
$S=(T-\lambda I)|_M$ is bounded below.  Otherwise there are $x_n\in M$ with
$\|x_n\|=1$ and $Sx_n\to0$.  Compactness gives a subsequence for which
$Tx_n$ converges, and
\[
  x_n=\lambda^{-1}(Tx_n-Sx_n)
\]
then converges to some $x\in M$ with $\|x\|=1$ and $Sx=0$, contradicting
$M\cap N=\{0\}$.  A bounded-below operator has closed range, and
$\Ran S=\Ran(T-\lambda I)$.
\end{proof}

\begin{notetheorem}[Stabilization and Riesz decomposition]
Let $T$ be compact and $\lambda\ne0$.  There is $m\ge1$ such that
\[
  \Ker(T-\lambda I)^m=\Ker(T-\lambda I)^{m+1}=\cdots,
\]
\[
  \Ran(T-\lambda I)^m=\Ran(T-\lambda I)^{m+1}=\cdots,
\]
and
\[
  X=\Ker(T-\lambda I)^m\oplus\Ran(T-\lambda I)^m.
\]
\end{notetheorem}

\begin{proof}
Write
\[
  S=T-\lambda I,
  \qquad
  N_n=\Ker S^n,
  \qquad
  R_n=\Ran S^n,
\]
with $N_0=\{0\}$ and $R_0=X$.  The spaces $N_n$ are closed.  The ranges are
closed as well: by the binomial formula,
\[
  S^n=(-\lambda)^nI+K_n,
\]
where $K_n$ is a finite sum of positive powers of the compact operator $T$ and
is therefore compact.  The earlier closed-range proposition applies to
$S^n$ because $(-\lambda)^n\ne0$.

\emph{Stabilization of the kernels.}
Suppose, contrary to the claim, that
$N_{n-1}\subsetneq N_n$ for every $n\ge1$.  Riesz's lemma gives
$x_n\in N_n$ such that
\[
  \|x_n\|=1,
  \qquad
  \dist(x_n,N_{n-1})>\frac12.
\]
If $n>m$, then $x_m\in N_m\subseteq N_{n-1}$,
$Sx_n\in N_{n-1}$, and $Tx_m=(S+\lambda I)x_m\in N_m\subseteq N_{n-1}$.
Consequently
\[
  Tx_n-Tx_m
  =\lambda x_n+Sx_n-Tx_m
  =\lambda x_n-z_{m,n}
\]
for some $z_{m,n}\in N_{n-1}$.  Hence
\[
  \|Tx_n-Tx_m\|
  =|\lambda|\left\|x_n-\frac{z_{m,n}}{\lambda}\right\|
  >\frac{|\lambda|}{2}.
\]
This contradicts compactness of $T$.  Thus $N_r=N_{r+1}$ for some $r$.
Once equality occurs it persists: if $N_r=N_{r+1}$ and
$x\in N_{r+2}$, then $Sx\in N_{r+1}=N_r$, so
$x\in N_{r+1}=N_r$; induction gives $N_r=N_{r+k}$ for all $k\ge1$.

\emph{Stabilization of the ranges.}
Suppose that $R_n\supsetneq R_{n+1}$ for every $n\ge0$.  Since the $R_n$ are
closed, Riesz's lemma gives $x_n\in R_n$ such that
\[
  \|x_n\|=1,
  \qquad
  \dist(x_n,R_{n+1})>\frac12.
\]
For $n>m$, the descending-chain relation gives
$x_n\in R_n\subseteq R_{m+1}$.  Also $Sx_m\in R_{m+1}$ and
$Tx_n=(S+\lambda I)x_n\in R_n\subseteq R_{m+1}$.  Therefore
\[
  Tx_m-Tx_n=\lambda x_m-z_{m,n}
  \qquad(z_{m,n}\in R_{m+1}),
\]
so
\[
  \|Tx_m-Tx_n\|>\frac{|\lambda|}{2},
\]
again contradicting compactness.  Hence $R_s=R_{s+1}$ for some $s$, and then
$R_s=R_{s+k}$ for every $k\ge1$.

This is the null-space/range argument appearing in the handwritten proof.  To
make its final step explicit, choose
\[
  m\ge\max\{r,s,1\}.
\]
Then $N_m=N_{2m}$ and $R_m=R_{2m}$.  Given $x\in X$, the vector $S^mx$ lies
in $R_m=R_{2m}$, so there is $z\in X$ with
\[
  S^mx=S^{2m}z.
\]
Thus
\[
  x=(x-S^mz)+S^mz,
\]
where $x-S^mz\in N_m$ and $S^mz\in R_m$.  Therefore $X=N_m+R_m$.
If $x\in N_m\cap R_m$, write $x=S^my$.  Then
$0=S^mx=S^{2m}y$, so $y\in N_{2m}=N_m$ and hence $x=S^my=0$.  The sum is
direct, proving
\[
  X=N_m\oplus R_m.
\]
\end{proof}

\section{Fredholm alternative}

\begin{notelemma}[Biorthogonal vectors]
If $f_1,\ldots,f_n\in X^*$ are linearly independent, then there exist
$x_1,\ldots,x_n\in X$ such that
\[
  f_i(x_j)=\delta_{ij}.
\]
\end{notelemma}

\begin{proof}
The map $Q\colon X\to\mathbb K^n$, $Qx=(f_1(x),\ldots,f_n(x))$, is
surjective; otherwise a nonzero functional on $\mathbb K^n$ annihilating
$Q(X)$ would give a nontrivial linear relation among the $f_i$.  Choose
$x_j$ with $Qx_j=e_j$.
\end{proof}

\begin{notetheorem}[Fredholm alternative]
Let $T\in\Kop{X}{X}$ and let $\lambda\ne0$.  Put
$S=\lambda I-T$.  Then:
\begin{enumerate}
  \item $S$ is injective if and only if it is surjective; in this case
        $S^{-1}$ is bounded;
  \item $\dim\Ker S=\dim\Ker S^*<\infty$;
  \item the equation $Sx=y$ is solvable if and only if
        \[
          f(y)=0\qquad(f\in\Ker S^*);
        \]
  \item the equation $S^*f=g$ is solvable if and only if
        \[
          g(x)=0\qquad(x\in\Ker S).
        \]
\end{enumerate}
Here $S^*=\lambda I-T^*$ for the Banach-space adjoint defined on linear
functionals.  Under the Hilbert-space adjoint identification, the corresponding
spectral parameter is conjugated according to the inner-product convention.
\end{notetheorem}

\begin{proof}
Since $S=\lambda I-T$ differs only by a sign from $T-\lambda I$, the
closed-range and stabilization results above apply to $S$ and all its powers.

We first follow the null-space/range argument of the notes.  If $S$ is
injective and $\Ran S\ne X$, then every inclusion
\[
  X\supseteq\Ran S\supseteq\Ran S^2\supseteq\cdots
\]
is strict.  Indeed, if $\Ran S^n=\Ran S^{n+1}$ for some $n\ge1$ and
$y\in\Ran S^{n-1}$, then $Sy\in\Ran S^n=\Ran S^{n+1}$, so
$Sy=S(S^nz)$ for some $z$; injectivity gives $y=S^nz\in\Ran S^n$.
Repeating this backward would yield $X=\Ran S$, a contradiction.  The strict
chain contradicts range stabilization.  Hence injectivity implies
surjectivity.

Conversely, suppose $S$ is surjective but has $0\ne x_1\in\Ker S$.  Choose
recursively $x_{n+1}$ with $Sx_{n+1}=x_n$.  Then
\[
  x_n\in\Ker S^n\setminus\Ker S^{n-1},
\]
so the kernels increase strictly forever, contradicting kernel stabilization.
Thus surjectivity implies injectivity.  If either condition holds, $S$ is a
bounded bijection of Banach spaces and the bounded inverse theorem gives
$S^{-1}\in\Bop{X}{X}$.

We next compute the dimensions.  Choose $m$ from the Riesz decomposition and
write
\[
  X=N_m\oplus R_m,
  \qquad N_m=\Ker S^m,
  \qquad R_m=\Ran S^m.
\]
The space $N_m$ is finite dimensional, because
$S^m=cI+K$ with $c\ne0$ and $K$ compact.  Since the ranges have stabilized,
$S(R_m)=R_m$; the direct-sum relation also shows that $S|_{R_m}$ is injective.
Moreover,
\[
  \Ran S=S(N_m)\oplus R_m.
\]
Consequently,
\[
\begin{aligned}
  \operatorname{codim}\Ran S
   &=\dim N_m-\dim S(N_m)\\
   &=\dim\Ker(S|_{N_m})
    =\dim\Ker S.
\end{aligned}
\]
Because $\Ran S$ is closed,
\[
  \Ker S^*=(\Ran S)^\perp
\]
is naturally the dual of the finite-dimensional quotient
$X/\Ran S$.  Therefore
\[
  \dim\Ker S^*=\operatorname{codim}\Ran S=\dim\Ker S<\infty.
\]

Finally, the standard annihilator identities give
\[
  \Ran S=\preann{\Ker S^*}
\]
because $\Ran S$ is closed.  Hence $Sx=y$ is solvable exactly when
$f(y)=0$ for every $f\in\Ker S^*$.  By Schauder's theorem, $T^*$ is compact,
so $S^*=\lambda I-T^*$ also has closed range.  The dual annihilator identity
then gives
\[
  \Ran S^*=(\Ker S)^\perp,
\]
which is precisely the stated solvability criterion for $S^*f=g$.
\end{proof}

\begin{notecorollary}
Every nonzero spectral value of a compact operator is an eigenvalue.
\end{notecorollary}

\begin{proof}
If $\lambda\ne0$ and $\lambda I-T$ were injective, the Fredholm alternative
would make it surjective and hence invertible.  Thus a nonzero point of the
spectrum must have nontrivial kernel.
\end{proof}

\section{Arzelà--Ascoli and integral operators}

\begin{notetheorem}[Arzelà--Ascoli theorem]
Let $K$ be a compact metric space.  A bounded equicontinuous family
$\mathcal F\subseteq C(K)$ is relatively compact in the supremum norm.
Equivalently, every sequence in $\mathcal F$ has a uniformly convergent
subsequence.
\end{notetheorem}

\begin{notetheorem}[Continuous-kernel integral operators]
Let $J=[a,b]$ and let $k\in C(J\times J)$.  Define
\[
  (Tf)(s)=\int_a^b k(s,t)f(t)\,\dd t,
  \qquad f\in C(J).
\]
Then $T\colon C(J)\to C(J)$ is compact.
\end{notetheorem}

\begin{proof}
For $\|f\|_\infty\le1$,
\[
  \|Tf\|_\infty
  \le(b-a)\|k\|_\infty,
\]
so the image of the unit ball is uniformly bounded.  Uniform continuity of
$k$ gives
\[
 |Tf(s_1)-Tf(s_2)|
 \le\int_a^b|k(s_1,t)-k(s_2,t)|\,\dd t,
\]
which tends to zero uniformly in $f$ as $s_1\to s_2$.  Thus the image is
equicontinuous, and Arzelà--Ascoli gives relative compactness.
\end{proof}
